% =====================================================================
%  Le protocole TRIADE — édition académique v0.4 (article de recherche)
%  Corps du document : compilation via pdf.py convert.latex (Tectonic).
%  La couverture est générée séparément (cover.html, Template 03) et
%  fusionnée en amont du fichier final — AUCUNE page de titre ici.
% =====================================================================
\documentclass[10pt,a4paper,twocolumn]{article}

% ---------- Polices (système, via fontspec/XeTeX) ----------
\usepackage{fontspec}
\setmainfont{Liberation Serif}          % texte — métrique Times
\setsansfont{Liberation Sans}           % titres de section
\setmonofont{Liberation Mono}           % code
\usepackage[french]{babel}
\addto\captionsfrench{\renewcommand{\tablename}{Tableau}}

% ---------- Mathématiques ----------
\usepackage{amsmath,amssymb,amsthm}

% ---------- Mise en page ----------
\usepackage{geometry}
\geometry{a4paper, top=2.3cm, bottom=2.2cm, left=1.9cm, right=1.9cm, headheight=14pt}
\usepackage{graphicx}
\usepackage{xcolor}
\definecolor{liaison}{HTML}{3681A6}     % accent série TRIADE (froid)
\definecolor{grisdoux}{HTML}{555F66}
\usepackage{booktabs}
\usepackage{tabularx}
\newcolumntype{Y}{>{\raggedright\arraybackslash}X}
\usepackage{multirow}
\usepackage{adjustbox}
\usepackage{enumitem}
\setlist{noitemsep, topsep=2pt, leftmargin=1.4em}
\usepackage{microtype}

% ---------- Théorèmes ----------
\theoremstyle{definition}
\newtheorem{defn}{Définition}[section]
\newtheorem{invariant}[defn]{Invariant}
\theoremstyle{plain}
\newtheorem{prop}[defn]{Proposition}

% ---------- Algorithmes ----------
\usepackage[ruled,vlined,linesnumbered]{algorithm2e}
\SetAlFnt{\small}
\SetAlCapFnt{\small}
\SetAlCapNameFnt{\small}
\SetAlgorithmName{Algorithme}{algorithme}{Liste des algorithmes}
\SetKwInput{KwIn}{Entrées}
\SetKwInput{KwOut}{Sorties}
\SetKw{KwTo}{à}

% ---------- Listings (code C#) ----------
\usepackage{listings}
\definecolor{codekeyword}{HTML}{1F3A54}
\definecolor{codecomment}{HTML}{5A7055}
\definecolor{codestring}{HTML}{7A4A1F}
\lstdefinelanguage{cs3}{
  morekeywords={public,const,double,var,new,if,throw,return,foreach,for,in,
    static,string,char,bool,void,class,private,readonly,int,float,using,
    namespace,partial,summary},
  sensitive=true,
  morecomment=[l]{//},
  morecomment=[l]{///},
  morestring=[b]",
}
\lstset{
  language=cs3,
  basicstyle=\footnotesize\ttfamily,
  keywordstyle=\bfseries\color{codekeyword},
  commentstyle=\itshape\color{codecomment},
  stringstyle=\color{codestring},
  numbers=none,
  breaklines=true,
  breakatwhitespace=true,
  postbreak=\mbox{\textcolor{grisdoux}{$\hookrightarrow$}\space},
  showstringspaces=false,
  columns=fullflexible,
  keepspaces=true,
  frame=single,
  framerule=0.3pt,
  rulecolor=\color{black!25},
  xleftmargin=4pt, xrightmargin=4pt,
  aboveskip=6pt, belowskip=6pt,
}

% ---------- Bibliographie ----------
\usepackage[numbers,sort&compress]{natbib}
\bibliographystyle{unsrtnat}
\setlength{\bibsep}{2pt}

% ---------- TikZ ----------
\usepackage{tikz}
\usetikzlibrary{arrows.meta,positioning}

% ---------- En-têtes / pieds ----------
\usepackage{fancyhdr}
\pagestyle{fancy}
\fancyhf{}
\fancyhead[L]{\small\itshape Le protocole TRIADE — édition académique v0.4}
\fancyhead[R]{\small\itshape}
\fancyfoot[C]{\small\thepage}
\renewcommand{\headrulewidth}{0.3pt}

% ---------- Divers ----------
\raggedbottom
\widowpenalty=10000
\clubpenalty=10000
\usepackage{titlesec}
\titleformat{\section}
  {\large\sffamily\bfseries\color{liaison}}{\thesection}{0.6em}{}
\titleformat{\subsection}
  {\normalsize\sffamily\bfseries}{\thesubsection}{0.6em}{}
\titlespacing*{\section}{0pt}{12pt plus 4pt minus 2pt}{6pt plus 2pt minus 2pt}

% hyperref : TOUJOURS en dernier
\usepackage[colorlinks=true, linkcolor=liaison, citecolor=grisdoux,
            urlcolor=liaison, bookmarks=true, bookmarksnumbered=true,
            unicode=true]{hyperref}

\hypersetup{
  pdftitle={Le protocole TRIADE : gouvernance vérifiable du développement
            logiciel assisté par IA — édition académique v0.4},
  pdfauthor={Alexis FR-LOPEZ ; Agent Z.ai (GLM)},
  pdfsubject={Gouvernance du développement logiciel assisté par IA ;
              étude de cas RoomEq},
  pdfkeywords={développement assisté par IA, dette technique, conception
               par contrat, gouvernance logicielle, étude de cas}
}

% ---------- Raccourcis ----------
\newcommand{\IB}{\ensuremath{\mathcal{I}}}
\newcommand{\BB}{\ensuremath{\mathcal{B}}}
\newcommand{\CB}{\ensuremath{\mathcal{C}}}
\newcommand{\DB}{\ensuremath{\mathcal{D}}}
\newcommand{\lexsucc}{\ensuremath{\succ_{\mathrm{lex}}}}
\newcommand{\lexsucceq}{\ensuremath{\succeq_{\mathrm{lex}}}}
\newcommand{\code}[1]{\texttt{#1}}

\begin{document}

% =====================================================================
%  BLOC DE TITRE (pleine largeur) — aucun \maketitle
% =====================================================================
\twocolumn[{%
\begin{center}
  {\footnotesize\sffamily\color{grisdoux}\MakeUppercase{Article de recherche
   \,·\, Édition académique du protocole \textsc{Triade} \,·\, Document v0.4
   (noyau normatif v0.3)}}\\[10pt]
  {\LARGE\bfseries Le protocole TRIADE : gouvernance vérifiable\\[3pt]
   du développement logiciel assisté par IA}\\[9pt]
  {\normalsize Alexis FR-LOPEZ\textsuperscript{1}
   \,·\, Agent Z.ai (GLM)\textsuperscript{1}}\\[3pt]
  {\small\itshape\textsuperscript{1}Projet TRIADE — développement agent-first,
   dépôts privés du développeur}\\[8pt]
  {\small 26 octobre 2026}
\end{center}
\vspace{2pt}
\begin{center}
\begin{minipage}{0.88\textwidth}
  \small
  \textbf{Résumé} — Le développement assisté par modèles de langage déplace le
  problème de la qualité du code vers celui de la gouvernance du processus :
  l’agent qui écrit le code n’est pas l’agent qui le relira dans six mois, et
  rien ne garantit naturellement qu’une évolution respecte les intentions
  passées. Nous présentons TRIADE, un protocole qui contraint toute évolution
  d’un projet logiciel à s’exprimer à travers trois artefacts versionnés —
  spécifications d’intention diffables (couche I), artefact exécutable scellé
  (couche B), carte indexée et signée (couche C) — unis par un verrou croisé
  cryptographique, dix règles à sanction mécanique et une fonction objectif
  lexicographique. Nous formalisons ces concepts : invariants de cohérence,
  règles de transformation, oracle différentiel dérivé des contrats. Nous
  rapportons une étude de cas monocasale (N = 1) : le recodage complet d’une
  application audio temps réel (RoomEq, environ 15~kloc C\#/.NET~8) depuis une
  base scellée v2.2.2, en sept vagues gouvernées par le protocole —
  44 spécifications écrites avant le code, 42 modules cartographiés, 437 tests
  contre 228, zéro avertissement de compilation, parité fonctionnelle de 14
  familles sur 14 et empreintes reproductibles. Trois comparatifs de code
  avant/après illustrent la conversion systématique de la classe de défauts
  « contrat non écrit » en contrats nommés, testés et journalisés. Les menaces
  sur la validité sont traitées explicitement : cas unique sans condition
  témoin, auto-vérification par l’agent auteur, et écart conceptuel entre la
  refactorisation source démontrée et la génération binaire native promise par
  la spécification, dont la validation reste ouverte (jalons M0–M7).
  \par\vspace{5pt}
  \textbf{Mots-clés} : développement assisté par IA~; dette technique~;
  conception par contrat~; gouvernance logicielle~; vérification~;
  étude de cas.
\end{minipage}
\end{center}
\vspace{16pt}
}]

% =====================================================================
%  1. INTRODUCTION
% =====================================================================
\section{Introduction}\label{sec:intro}

Le développement logiciel est entré dans l’ère \emph{agent-first} : une part
croissante du code est écrite, relue et modifiée par des modèles de langage
plutôt que par des humains. Les études empiriques convergent sur l’ampleur du
phénomène (gains de productivité mesurés ou déclarés chez les utilisateurs
d’assistants de complétion \cite{ziegler2022,peng2023}, adoption massive et
réorganisation des pratiques de lecture et de vérification du code généré
\cite{vaithilingam2022,barke2023}) tandis que les revues systématiques
recensent déjà plusieurs centaines de travaux sur les modèles de langage
appliqués au génie logiciel \cite{hou2024,fan2023}. Cette bascule remet en
question une hypothèse implicite de nos outils : que le code source est optimisé
pour la \emph{lecture humaine}. Compilateurs, systèmes de construction et
langages ont été conçus autour de cette hypothèse ; si le lecteur principal
devient une machine, l’empilement mérite un réexamen systématique plutôt qu’une
révision marginale.

Or la question qui se pose au praticien n’est pas celle que mesurent la
plupart de ces travaux. La littérature évalue des \emph{capacités} (le modèle
peut-il produire un code correct pour une tâche donnée~?
\cite{chen2021,austin2021}) ou des \emph{effets} (la productivité change-t-elle~?
\cite{peng2023}). Elle documente aussi les modes d’échec : code plausible mais
inexact, dépendances hallucinées, conformité superficielle à la demande
\cite{hou2024}. Mais entre le modèle et le produit livré s’étend un territoire
peu exploré~: le \emph{processus} — l’ensemble des règles qui déterminent ce
qu’un agent a le droit d’écrire, quand une écriture devient permanente, et
comment un projet survit au changement d’agent. C’est ce territoire que
TRIADE se propose de gouverner.

TRIADE (\emph{Intention, Binaire, Carte}) est un protocole de gouvernance du
développement logiciel conçu pour un monde où l’exécutant est un agent. Son
principe est simple à énoncer~: \emph{aucune écriture n’est permanente sans
vérification, et aucune évolution ne commence sans intention écrite}. Toute
modification traverse trois artefacts versionnés — la couche I
(spécifications d’intention diffables), la couche B (l’artefact exécutable
scellé), la couche C (une carte indexée, signée, qui relie les deux) — et dix
règles à sanction mécanique (refus de construction, refus de commit, blocage
de session) plutôt qu’à recommandation. Un verrou croisé cryptographique relie
la carte et le binaire dans les deux sens ; une fonction objectif
lexicographique (sécurité $>$ performance $>$ optimisation $>$ coût en tokens)
arbitre tout conflit sans compensation possible entre rangs. Le protocole
existe en version~0.3, spécifié dans le dépôt du projet avec son journal de
décisions.

Cet article restructure cette spécification selon les standards d’un article de
recherche en génie logiciel et la confronte à sa première validation à
l’échelle~: le recodage complet de RoomEq, un égaliseur acoustique temps réel
d’environ 15~kloc C\#/.NET~8, exécuté en sept vagues depuis une base scellée
v2.2.2. Nos contributions sont les suivantes~:
\begin{itemize}
\item \textbf{C1 — Formalisation.} Nous donnons au protocole une présentation
formalisée~: état du projet comme quadruplet, prédicat de cohérence du verrou
croisé, règles de transformation en système de transitions, oracle
différentiel, fonction objectif lexicographique
(section~\ref{sec:methode}).
\item \textbf{C2 — Étude de cas à l’échelle.} Nous rapportons le recodage
complet d’une application réelle sous gouvernance TRIADE, avec ses métriques
(structure, défense, gouvernance), ses comparatifs de code avant/après et sa
séquence de publication — y compris deux échecs d’intégration continue honorés
puis corrigés (section~\ref{sec:etude}).
\item \textbf{C3 — Reverse-specs.} Nous introduisons la technique de
spécification rétro-active de binaires scellés hérités, qui étend la couche I
au patrimoine existant sans le recoder (section~\ref{sec:reverse}).
\item \textbf{C4 — Auto-évaluation critique.} Nous délimitons explicitement ce
que la démonstration établit et ne prouve pas, selon le cadre des menaces sur
la validité de Wohlin et al. \cite{wohlin2012} — notamment la validité externe
d’un cas unique et l’écart entre refactorisation source et génération binaire
native (section~\ref{sec:menaces}).
\end{itemize}

L’article suit l’ordre attendu~: la section~\ref{sec:connexes} situe TRIADE
par rapport à la dette technique, à la conception par contrat et au génie
logiciel assisté par IA~; la section~\ref{sec:methode} formalise le
protocole~; la section~\ref{sec:etude} présente l’étude de cas~; les
sections~\ref{sec:discussion} et~\ref{sec:menaces} discutent les résultats et
leurs limites~; la section~\ref{sec:conclusion} conclut. Le lecteur pressé
trouvera en tableau~\ref{tab:regles} la traduction de chaque règle dans le
dépôt de l’étude de cas — c’est la table qui a gouverné sept vagues de
recodage.

% =====================================================================
%  2. TRAVAUX CONNEXES
% =====================================================================
\section{Travaux connexes}\label{sec:connexes}

TRIADE emprunte à trois traditions~: la gestion de la dette technique, la
conception par contrat et, plus récemment, le génie logiciel assisté par IA
(AISE). Cette section les passe en revue successivement, puis positionne le
protocole par rapport à chacune.

\subsection{Dette technique}\label{sec:connexes-dette}

La métaphore de la dette technique, introduite par Cunningham, décrit la
conséquence d’un choix de conception opportun~: une facilité immédiate qui se
rembourse, avec intérêts, dans la maintenance future \cite{cunningham1992}.
Kruchten et al. en ont fait un objet de recherche à part entière, en
identifiant notamment la dimension \emph{sociale} et documentaire de la dette
— la perte de contexte des décisions passées \cite{kruchten2012}. La revue
systématique de Li et al. classe les activités de gestion en trois
familles~: identification, mesure et remboursement \cite{li2015}, et des
méthodes comme SQALE proposent d’évaluer le montant de la dette sur le code
existant \cite{letouzey2012}. Un courant connexe formalise la mémoire des
décisions elles-mêmes — les \emph{architecture decision records} de Nygard
\cite{nygard2011} — en réaction à la même perte de contexte.

Ces approches partagent un trait~: elles interviennent \emph{a posteriori},
sur un code dont la dette est déjà contractée. TRIADE inverse la perspective~:
la dette est traitée à la source, en rendant son principal explicite avant
l’écriture. La couche I est précisément l’artefact qui empêche la classe la
plus insidieuse de dette — le \emph{contrat non écrit} (valeur-seuil enfouie,
convention d’interface implicite, comportement dépendant de l’hôte) — d’être
contractée en silence~: ce qui n’est pas spécifié n’est pas codé, et ce qui
est codé renvoie à sa spécification. Le journal des décisions, requis par la
gouvernance du protocole, généralise les enregistrements de Nygard en les
rattachant à une fonction objectif explicite.

\subsection{Conception par contrat et spécification}\label{sec:connexes-dbc}

La lignée remonte à la logique axiomatique de Hoare, qui attache des
pré-conditions et post-conditions aux programmes \cite{hoare1969}, et à la
discipline de la programmation vérifiée de Dijkstra \cite{dijkstra1976}.
Meyer a popularisé la \emph{conception par contrat}~: chaque routine publie
ses obligations mutuelles — pré-conditions, post-conditions, invariants —
vérifiées à l’exécution \cite{meyer1992}~; Findler et Felleisen ont étendu la
notion aux fonctions d’ordre supérieur avec un contrôle dynamique des
violations \cite{findler2002}. Sur le plan de la vérification, QuickCheck a
établi que des propriétés déclaratives et des tests aléatoires contraints
formaient un oracle bon marché et redoutable \cite{claessen2000}. Pour le
code hérité, Feathers définit les \emph{characterization tests}~: figer le
comportement existant sous des tests avant d’y toucher \cite{feathers2004},
dans la tradition du remaniement contrôlé \cite{fowler1999}. Enfin, les
méthodes formelles légères (Alloy \cite{jackson2006}) ou lourdes (TLA+
\cite{lamport2002}) offrent des langages de spécification vérifiables par
model checking.

TRIADE se distingue des deux manières. D’une part, ses contrats sont
\emph{légers par décision explicite} (D3 du journal du projet)~: signatures
typées, pré/post-conditions en prose contrainte et invariants — pas de
sémantique formelle exécutable, au prix d’une vérification par tests plutôt
que par preuve. D’autre part et surtout, le contrat change de plan~: dans la
conception par contrat classique, le contrat est une annotation \emph{dans} le
code, surveillée à l’exécution~; dans TRIADE, le contrat est un artefact
\emph{extérieur} au code — diffable, versionné, antérieur à l’écriture — qui
\emph{conditionne le droit même de commiter} (règle R5). Le contrat n’est
plus une assertion runtime mais un instrument de gouvernance du processus.

\subsection{Génie logiciel assisté par IA}\label{sec:connexes-aise}

L’évaluation des modèles de code a d’abord mesuré des capacitudes~: réussite
sur des bancs d’essai comme HumanEval \cite{chen2021} ou la synthèse de
programmes à partir de spécifications en langage naturel \cite{austin2021}.
Un second courant mesure les effets sur les développeurs~: évaluations de
productivité sur des tâches contrôlées \cite{ziegler2022}, essai randomisé en
entreprise \cite{peng2023}, études d’utilisabilité mettant en évidence
l’écart entre attentes et expérience \cite{vaithilingam2022}, et analyses
fine-grained des modes d’interaction \cite{barke2023}. Les revues
systématiques agrègent ce corpus et y relèvent des angles morts récurrents~:
mesures de qualité au-delà des bancs d’essai, maintenance à long terme,
et confiance \cite{hou2024,fan2023}. Enfin, l’intégrité de la chaîne
d’approvisionnement a produit ses propres instruments~: constructions
reproductibles bit à bit \cite{lamb2021} et métadonnées de provenance signées
(in-toto) \cite{torresarias2019}.

À notre connaissance, la littérature AISE ne propose pas de protocole de
gouvernance du \emph{processus de développement} lui-même. Les travaux
existentants supposent l’outillage standard (répertoire versionné, revue de
code, intégration continue) et évaluent l’agent dans ce cadre hérité de
l’ère humaine. TRIADE suppose au contraire l’agent acquis et redessine le
cadre~: l’agent n’accède à l’artefact que par la carte (économie de
contexte), n’écrit qu’après vérification (R5), et voit ses évolutions
arbitrées par une fonction objectif explicite. Les verrouillages
cryptographiques s’inspirent des approches de provenance
\cite{torresarias2019,lamb2021}, mais protègent une relation inédite~: celle
entre l’\emph{index de développement} privé et l’artefact livré.

\subsection{Positionnement}\label{sec:connexes-position}

Le tableau~\ref{tab:position} résume le positionnement. Les approches de
dette technique mesurent un état et recommandent~; la conception par contrat
instrumente le code~; l’AISE évalue des modèles et des effets. TRIADE est le
seul des quatre à définir un \emph{protocole de processus}~: un ensemble
d’artefacts obligatoires, de règles à sanction mécanique et d’un ordre de
priorité non compensatoire, applicable à tout exécutant — humain ou agent.
Cette généralité est aussi sa limite~: rien, dans l’état actuel, ne prouve
que le protocole survive au contact d’équipes humaines nombreuses, dont les
incitations diffèrent de celles d’un agent (section~\ref{sec:menaces}).

\begin{table}[t]
\centering\small
\caption{Positionnement de TRIADE parmi les approches connexes.}
\label{tab:position}
\begin{tabularx}{\columnwidth}{@{}l Y Y@{}}
\toprule
Approche & Moment & Instrument pivot \\
\midrule
Dette technique \cite{li2015} & après coup & mesures, tableaux \\
Contrat (DbC) \cite{meyer1992} & à l’exécution & assertions dans le code \\
AISE \cite{hou2024} & évaluation & modèles, bancs d’essai \\
\textbf{TRIADE} (cet art.) & \textbf{avant écriture} & \textbf{specs + règles
à sanction} \\
\bottomrule
\end{tabularx}
\end{table}

% =====================================================================
%  3. MÉTHODOLOGIE — LE PROTOCOLE TRIADE
% =====================================================================
\section{Le protocole TRIADE}\label{sec:methode}

Cette section formalise le protocole tel que spécifié en version~0.3. La
présentation suit l’ordre des dépendances~: artefacts et état
(section~\ref{sec:artefacts}), fonction objectif
(section~\ref{sec:objectif}), invariants de cohérence
(section~\ref{sec:verrou}), règles et boucle de mise à jour
(section~\ref{sec:regles}), stratégie de vérification
(section~\ref{sec:verification}), extension aux cibles multiples et
gouvernance (section~\ref{sec:profils}).

\subsection{Artefacts et état du projet}\label{sec:artefacts}

\begin{defn}[État du projet]\label{def:etat}
Un projet TRIADE à l’instant $t$ est un quadruplet
\begin{equation}\label{eq:etat}
S_t \;=\; \langle\, \IB_t,\ \BB_t,\ \CB_t,\ \DB_t \,\rangle
\end{equation}
où $\IB_t$ est un ensemble fini de \emph{spécifications d’intention}
$\sigma = (\mathit{nom},\ \mathit{pre},\ \mathit{post},\ \mathit{inv},\
\mathit{adr})$~; $\BB_t$ l’artefact exécutable, décomposé en blocs
versionnés~; $\CB_t$ la \emph{carte}, index signé liant chaque adresse à ses
spécifications, empreintes et preuves~; $\DB_t$ le \emph{journal des
décisions}, append-only. Le couple $(\IB, \BB)$ est régi par deux
principes~: $\IB$ est écrit \emph{avant} $\BB$ (règle R2), et $\BB$ n’est
jamais livré sans $\IB$ et $\CB$ à jour (règle R1).
\end{defn}

\begin{defn}[Adresse stable]\label{def:adresse}
Soit $A$ l’espace des adresses. Une adresse $a \in A$ est \emph{stable} si
sa dénotation — le point d’entrée qu’elle désigne dans $\BB$ — ne change pas
d’une révision à l’autre (règle R7). La carte maintient une relation
d’indexation $\CB_{\mathrm{loc}} \subseteq A \times \BB$ telle que chaque
bloc soit décrit (règle R3) et chaque spécification indexée. La stabilité
des adresses est la condition de possibilité du \emph{patch à contrat
stable}~: remplacer le bloc dénoté par $a$ sans invalider aucun appelant.
\end{defn}

La figure~\ref{fig:arch} montre l’architecture~: le diff d’intention
descend de $\IB$ vers $\CB$, le verrou croisé relie $\CB$ et $\BB$ dans les
deux sens, et seule la couche $\BB$ quitte le laboratoire. Aucune couche
n’est optionnelle~: privé de $\IB$, le système perd son ancrage
sémantique~; privé de $\CB$, l’agent devient aveugle~; privé de $\BB$, il
n’y a plus rien à exécuter.

\begin{figure}[t]
\centering
\resizebox{0.86\columnwidth}{!}{%
\begin{tikzpicture}[
  box/.style={draw, rounded corners=2pt, text width=4.6cm,
              minimum height=1.05cm, align=center,
              font=\small\sffamily},
  arr/.style={-{Stealth[length=5pt]}, thick, gray!65},
  note/.style={font=\scriptsize\sffamily, text=gray!80!black}
]
  \node[box, fill=blue!7]  (i) at (0,3.0)
    {\textbf{Couche I — Intention}\\spécifications diffables :\\contrats, pré/post, invariants};
  \node[box, fill=green!8] (c) at (0,1.5)
    {\textbf{Couche C — Carte}\\index signé, chiffré, fragmenté :\\adresses, preuves, manifestes};
  \node[box, fill=orange!9] (b) at (0,0)
    {\textbf{Couche B — Binaire}\\blocs scellés, adresses stables,\\checksums au chargement};
  \draw[arr] (i) -- node[right, note] {diff d’intention} (c);
  \draw[arr] (c) -- node[right, note] {localisation, régénération} (b);
  \draw[arr] (b.west) .. controls (-2.5,0.9) and (-2.5,2.1) .. node[left, note] {empreintes} (c.west);
  \draw[arr] (c.east) .. controls (2.5,2.1) and (2.5,0.9) .. node[right, note] {signature, scellage} (b.east);
  \draw[arr] (b.south) -- (0,-0.85) node[below, note] {livraison (B seul)};
\end{tikzpicture}}
\caption{Architecture TRIADE. Le verrou croisé relie C et B dans les deux
sens~: empreintes des blocs portées par la carte (gauche), signature de la
carte et scellage du binaire (droite). Seule la couche B est livrée.}
\label{fig:arch}
\end{figure}

\subsection{Fonction objectif lexicographique}\label{sec:objectif}

Le protocole arbitre toute évolution par un ordre de priorité non
compensatoire~: sécurité $>$ performance $>$ optimisation $>$ coût en
tokens. La formalisation naturelle est l’ordre lexicographique.

\begin{defn}[Admissibilité d’une évolution]\label{def:lex}
Soit une évolution $e : S \to S'$. On lui associe le vecteur de variations
\begin{equation}\label{eq:delta}
\Delta(e) = \bigl(\Delta_{\mathrm{sec}},\ \Delta_{\mathrm{perf}},\
\Delta_{\mathrm{opt}},\ \Delta_{\mathrm{tok}}\bigr) \in \mathbb{R}^4
\end{equation}
où chaque composante quantifie l’effet de $e$ sur le critère correspondant,
le coût en tokens étant compté négativement (un coût moindre est un gain).
L’évolution $e$ est \emph{admissible} si et seulement si
\begin{equation}\label{eq:lex}
\mathit{adm}(e) \;\Longleftrightarrow\; \Delta(e)\ \lexsucceq\ \mathbf{0}.
\end{equation}
\end{defn}

L’ordre lexicographique $\lexsucceq$ compare les vecteurs composante par
composante dans l’ordre des priorités~: $\Delta \lexsucc \mathbf{0}$ dès que
la première composante non nulle est strictement positive. La propriété
décisive est l’\emph{absence de compensation}~: aucun gain, si grand soit-il,
sur un critère de rang inférieur ne rachète une perte, si petite soit-elle,
sur un critère de rang supérieur. Concrètement~: aucune économie de tokens
ne justifie de sauter une vérification (sécurité)~; aucun raffinement
d’optimisation ne justifie une couche d’abstraction qui dégraderait les
performances de l’artefact. La définition~\ref{def:lex} jouera un rôle
d’oracle de gouvernance dans l’étude de cas~: c’est elle qui a motivé le
refus, mesuré à l’appui, d’une demande de fusion de modules et de
réécriture « en natif machine » (section~\ref{sec:resultats}).

\subsection{Verrou croisé et invariants}\label{sec:verrou}

Le verrou croisé répond à une menace précise~: l’empoisonnement de la
carte. Un attaquant — ou une dérive de l’agent — qui falsifierait les
métadonnées ferait travailler l’agent sur des informations fausses, et
pourrait lui faire introduire une porte dérobée en croyant respecter les
contrats. La contre-mesure est bilatérale et se formalise par un prédicat
de cohérence.

\begin{defn}[Cohérence du verrou croisé]\label{def:kappa}
L’état $S = \langle \IB, \BB, \CB, \DB \rangle$ est \emph{cohérent}, noté
$\kappa(S) = 1$, si et seulement si les cinq conjoints suivants
tiennent~:
\begin{align}
&\text{(i)}\quad \forall \sigma \in \IB :\ \mathit{adr}(\sigma) \in
\mathrm{dom}\,\CB_{\mathrm{loc}}
\label{eq:k1}\\
&\text{(ii)}\quad \forall b \in \BB :\ H(b) \in \CB_{\mathrm{bld}}
\label{eq:k2}\\
&\text{(iii)}\quad \CB.h_{\BB} = H(\BB) \;\wedge\;
\CB.h_{\IB} = H(\IB)
\label{eq:k3}\\
&\text{(iv)}\quad \mathrm{Verify}\bigl(pk_{\mathrm{dev}},\ H(\CB),\
\CB.\sigma\bigr) = 1
\label{eq:k4}\\
&\text{(v)}\quad \mathcal{O}_{\IB}(\BB) = \top
\label{eq:k5}
\end{align}
c’est-à-dire~: toute spécification est indexée~\eqref{eq:k1}, tout bloc est
empreinté~\eqref{eq:k2}, les empreintes portées par la carte correspondent
aux artefacts réels~\eqref{eq:k3}, la carte est signée par la clé du
développeur~\eqref{eq:k4}, et l’oracle dérivé des contrats accepte
l’artefact~\eqref{eq:k5}.
\end{defn}

Quatre \emph{invariants structurels}, vérifiés à chaque cycle de
construction, en spécialisent la portée~:

\begin{invariant}[Jamais de binaire nu — R1]\label{inv:r1}
Tout artefact $\BB$ est accompagné de $\IB$ et $\CB$ à jour. Violation~:
refus de construction.
\end{invariant}

\begin{invariant}[Adresses stables — R7]\label{inv:r7}
Les points d’entrée exposés ne se déplacent pas~; un remplacement préserve
la dénotation de chaque adresse, par profil de cible. Violation~:
re-liage forcée, interdite.
\end{invariant}

\begin{invariant}[Carte signée — R9]\label{inv:r9}
L’agent refuse toute carte non signée~; la signature couvre la carte
entière, en-tête compris. Violation~: refus de démarrage de session.
\end{invariant}

\begin{invariant}[Échec contrôlé au chargement — R10]\label{inv:r10}
L’artefact vérifie ses propres checksums au lancement~; tout écart
provoque un arrêt propre. Violation~: arrêt de l’exécution.
\end{invariant}

La sécurité du dispositif ne repose jamais sur le secret de la carte —
principe de Kerckhoffs \cite{kerckhoffs1883} — mais sur la difficulté de
falsifier simultanément les deux directions du verrou. La conception
s’apparente aux métadonnées de provenance signées des chaînes
d’approvisionnement \cite{torresarias2019} et aux constructions
reproductibles \cite{lamb2021}, appliquées à une relation inédite~: celle
entre l’index de développement privé et l’artefact livré.

\subsection{Règles et boucle de mise à jour}\label{sec:regles}

Le protocole compte dix règles (R1–R10), héritées des versions 0.1 et 0.2
et inchangées depuis. Elles sont peu nombreuses, vérifiables
mécaniquement, et sanctionnées par des effets concrets — refus de
construction, refus de commit, blocage de session — plutôt que par des
recommandations. Le tableau~\ref{tab:regles} les énonce avec leur
traduction dans le dépôt de l’étude de cas
(section~\ref{sec:etude})~; la présente sous-section en donne la lecture
formelle.

\begin{table*}[t]
\centering\small
\caption{Les dix règles du protocole et leur instanciation dans le dépôt de
l’étude de cas RoomEq (section~\ref{sec:etude}).}
\label{tab:regles}
\begin{tabularx}{\textwidth}{@{}l l Y Y@{}}
\toprule
\textbf{\#} & \textbf{Règle} & \textbf{Énoncé (spécification)} &
\textbf{Traduction RoomEq} \\
\midrule
R1 & Jamais de binaire nu & Tout artefact B est accompagné de I et C à jour
& Aucun fichier de \code{src/} sans module cartographié ni spécification I \\
R2 & Intention d’abord & Toute modification commence par un diff sur I &
La spécification du module est écrite avant le code (44 spécifications) \\
R3 & Balisage obligatoire & Chaque bloc est décrit dans C & Chaque commit
porte son module~; la carte est tenue à jour \\
R4 & Patch à contrat stable & Le remplacement préserve signatures et
contrats~; les appelants ne sont jamais invalidés & L’API publique d’un
module ne change pas sans révision de sa spécification \\
R5 & Vérifie-puis-écrit & Aucun commit sans oracle, vecteurs dorés et
round-trip & Construction 0 erreur / 0 avertissement et 100~\% des tests
verts avant chaque commit \\
R6 & Cristallisation & Dix vérifications réussies rendent le bloc
immuable~; le modifier exige la procédure lourde & Balises de version~;
42 modules verrouillés au cristal final \\
R7 & Adresses stables & Les points d’entrée sont garantis par padding~;
patcher ne déplace rien & Identifiants de module \code{M-<DOM>-NN}
immuables dans la carte \\
R8 & Verrou croisé & Empreintes des blocs et signature de la carte
vérifiées avant chaque édition & Empreintes SHA-256 des données critiques~;
parité md5 vérifiée \\
R9 & Carte signée obligatoire & L’agent refuse toute carte non signée &
Historique git intégral~; base de référence scellée (\code{deff25c}) \\
R10 & Échec contrôlé au chargement & Checksums vérifiés au lancement~;
écart = arrêt propre & Tests d’échec explicites (fragments courts, bruit
seul, etc.) \\
\bottomrule
\end{tabularx}
\end{table*}

Formellement, les règles définissent un \emph{système de transitions
étiquetées}~: $S \xrightarrow{\;r_i\;} S'$ n’est licite que si la
pré-condition de $r_i$ est satisfaite dans $S$ et sa post-condition
établissable dans $S'$. Les deux règles les plus chargées sont R5 et R6.
R5 (vérifie-puis-écrit) fait de la vérification une \emph{garde de commit}~:
\begin{equation}\label{eq:r5}
\mathit{commit}(S') \;\Rightarrow\; \kappa(S') \,\wedge\,
\mathcal{O}_{\IB'}(\BB') = \top \,\wedge\, \mathit{build}(\BB') = (0,0)
\end{equation}
autrement dit, aucune écriture n’est permanente sans cohérence du verrou,
verdict favorable de l’oracle et construction propre. R6
(cristallisation) associe à chaque bloc un compteur de vérifications
réussies $v(b) \in \mathbb{N}$, croissant~; lorsque $v(b) \geq 10$, le bloc
devient immuable et toute modification exige une \emph{nouvelle version} —
l’ancienne restant archivée, jamais détruite en silence. Les balises de
version du dépôt jouent exactement ce rôle dans l’étude de cas, avec la
nuance des pré-versions \code{-dev.N}, inventées pendant la démonstration
pour distinguer « vérifié par l’oracle » de « validé sur machine
réelle ».

L’algorithme~\ref{alg:boucle} décrit la boucle de mise à jour — l’unique
voie par laquelle le logiciel évolue. Elle se déroule en sept étapes
réparties en deux phases, l’édition puis la validation et le commit,
séparées par un point de non-retour~: rien n’est écrit dans l’artefact
avant que l’oracle n’ait rendu son verdict.

\begin{algorithm}[t]
\SetAlgoLined
\KwIn{état cohérent $S = \langle \IB, \BB, \CB, \DB \rangle$, requête de
changement $r$}
\KwOut{état cohérent $S'$ incorporant $r$, ou refus motivé}
\tcp{Phase d’édition}
$\delta \leftarrow$ diff minimal sur $\IB$ satisfaisant $r$
\quad\textnormal{(R2)}\;
$X \leftarrow \CB_{\mathrm{loc}}.\mathit{localiser}(\delta)$
\quad\textnormal{(R3)}\;
\If{$\kappa(S) = 0$}{\textbf{refuser} — session bloquée
\textnormal{(R8, R9)}\;}
$\BB^{*} \leftarrow$ régénérer les blocs de $X$ sous contrats stables
\textnormal{(R4, R7)}\;
\tcp{Point de non-retour}
\tcp{Phase de validation}
\If{$\mathcal{O}_{\IB'}(\BB^{*}) = \top$ \textnormal{et}
$\mathit{build}(\BB^{*}) = (0,0)$}{
  $S' \leftarrow \mathit{commit}(\IB', \BB^{*}, \CB', \DB')$ — patch,
  carte, scellage atomiques \textnormal{(R1, R5)}\;
  $\forall b \in X : v(b) \leftarrow v(b) + 1$~;
  \If{$v(b) \geq 10$}{cristalliser $b$ \textnormal{(R6)}\;}
}
\Else{\textbf{refuser} — candidats rejetés, journal consigné\;}
\caption{Boucle de mise à jour TRIADE.}
\label{alg:boucle}
\end{algorithm}

Le coût en tokens — seul des quatre objectifs spécifique à l’ère agent-first
— est traité comme une contrainte d’ingénierie~: l’agent ne lit jamais
l’artefact en entier, la carte est son unique index, conçue pour être lue
parcellairement. Les budgets normatifs s’expriment en inégalités~:
localiser un bloc $\leq 40$ tokens, éditer simplement $\leq 3\,000$,
ajouter $\leq 6\,000$~; l’audit complet de la carte est une opération
exceptionnelle, non budgétée. Ces bornes sont des \emph{cibles de
conception} de la carte fragmentée, pas encore des mesures — point
repris dans les menaces sur la validité (section~\ref{sec:menaces}).

\subsection{Stratégie de vérification}\label{sec:verification}

Supprimer le compilateur du pipeline — l’objectif à long terme du
protocole — revient à perdre les filets de sécurité qu’il tendait
gratuitement~: vérification de types, initialisation, analyse de flot. La
stratégie de vérification reconstruit une confiance équivalente sans
réintroduire de couche de traduction dans l’artefact. Elle repose sur
trois piliers.

\paragraph{Oracle éphémère.}
L’oracle est une implémentation de référence de chaque bloc — triviale,
volontairement lente, écrite dans un langage sûr et \emph{dérivée de la
couche I}, générée au début de chaque session puis détruite après usage~:
jamais livrée, jamais archivée. Le test est différentiel par
propriétés~:
\begin{equation}\label{eq:oracle}
\forall x \models \mathit{pre}(\sigma) :\quad
\mathrm{Obs}\bigl(b(x)\bigr) = \mathrm{Obs}\bigl(\mathcal{O}_{\IB}(x)\bigr)
\end{equation}
où $x$ parcourt des entrées aléatoires contraintes par les pré-conditions
du contrat, à la manière de QuickCheck \cite{claessen2000}, avec un volume
de référence de douze mille cas par bloc critique. Trois propriétés
rendent l’oracle fiable~: il dérive de $\IB$ et non de $\BB$ (il ne peut
pas hériter d’un défaut du code)~; il est régénéré à chaque session (jamais
en désaccord avec les contrats courants)~; sa régénération coûte moins que
le diagnostic d’une corruption découverte après scellage. L’oracle étant
dérivé de la couche I — indépendante de la cible — le même oracle valide
les blocs de tous les profils~: la portabilité ne divise pas le coût de
vérification, elle le mutualise.

\paragraph{Round-trip de notation.}
Chaque bloc assemblé par le moteur de notation du profil est désassemblé
par le vérificateur correspondant, et les instructions obtenues sont
comparées à la notation d’origine. Tout désaccord — une instruction, un
opérande, une taille d’encodage — est un échec de construction. Ce contrôle
à coût quasi nul attrape la classe d’erreurs la plus insidieuse~: l’écart
silencieux entre ce que l’agent croit avoir écrit et les octets réellement
produits. La nuance capitale est que la notation n’est \emph{pas} un
compilateur — là où une chaîne moderne applique ses passes d’optimisation
entre la représentation d’entrée et le code machine \cite{lattner2004},
aucune réécriture ne s’intercale ici entre l’écriture et les octets
(principe WYSIWYG).

\paragraph{Cristallisation.}
Plus un bloc fait ses preuves, moins il coûte à vérifier~: le compteur
$v(b)$ déplace le budget de vérification vers les blocs jeunes et les
zones actives. La chaîne de confiance du scellage résume le parcours d’une
preuve~: verdict de l’oracle $\to$ vecteurs dorés hachés dans la carte
$\to$ empreinte du bloc $\to$ signature de la carte entière $\to$
immuabilité à dix vérifications.

\subsection{Profils de cible, gouvernance, continuité}\label{sec:profils}

Depuis la version 0.3, le cœur du protocole — couches, règles, boucle,
vérification — est indépendant de la cible par construction~: tout ce qui
dépend de la cible est regroupé dans un \emph{profil} défini par sept
éléments vérifiables (ISA et types, ABI, format d’exécutable et
protections, interface système, mécanisme de patch à chaud, moteur de
notation, granularité et padding). Instancier une plateforme nouvelle
revient à écrire un profil, jamais à modifier le protocole. Le
tableau~\ref{tab:profils} résume les cinq profils déclarés~; le profil
primaire (Windows~11, x86-64, PE32+) porte la feuille de route M0–M5, et
les jalons M6/M7 valideront l’universalité en isolant chaque variable —
changement d’OS à ISA constante, puis changement d’ISA.

\begin{table}[t]
\centering\small
\caption{Profils de cible déclarés en v0.3.}
\label{tab:profils}
\begin{tabularx}{\columnwidth}{@{}l l Y@{}}
\toprule
\textbf{Profil} & \textbf{Format} & \textbf{Statut} \\
\midrule
\code{x86-64-win11-pe} & PE32+ & primaire, jalons M0–M5 \\
\code{x86-64-linux-elf64} & ELF64 & M6 — inter-OS \\
\code{riscv64-linux-elf64} & ELF64 & M7 — inter-ISA \\
\code{aarch64-macos-macho} & Mach-O & réservé \\
\code{wasm32-wasi} & Wasm & réservé, étude \\
\bottomrule
\end{tabularx}
\end{table}

La gouvernance complète le dispositif. Le journal des décisions $\DB$
consigne toute décision structurante avec ses alternatives rejetées et sa
justification rattachée à la fonction objectif~; une révision doit soit
respecter les décisions passées, soit expliciter précisément ce qu’elle
invalide — rien ne se change en silence, dans l’esprit des enregistrements
de décisions d’architecture \cite{nygard2011} étendus à un registre
normatif. Enfin, l’\emph{export de continuité} garantit l’anti-enfermement~:
un conteneur unique, chiffré et signé, regroupe les couches $\IB$, $\BB$,
$\CB$ complètes et un manifeste autodescriptif~; l’oracle, éphémère par
conception, y est volontairement exclu et régénéré depuis $\IB$ à
l’arrivée. La procédure d’import en cinq étapes — déchiffrement,
vérification de signature, vérification du verrou croisé, régénération de
l’oracle, reprise — applique au transfert les mêmes garde-fous qu’à
l’édition~: rien n’est interprété avant vérification, et l’échec d’une
étape est terminal.

% =====================================================================
%  4. ÉTUDE DE CAS — RECODAGE DE ROOMEQ
% =====================================================================
\section{Étude de cas : recodage de RoomEq}\label{sec:etude}

\subsection{Contexte et protocole expérimental}\label{sec:protocole-etude}

RoomEq est un égaliseur acoustique temps réel pour Windows et macOS~: il
mesure la réponse d’une pièce par balayage logarithmique, calcule une
correction, l’applique en temps réel sur le son du système, et offre une
console de mixage complète — pistes, égaliseur par piste, routage par carte
son virtuelle, pont ASIO vers les stations de travail audio. La base de
départ est la version 2.2.2, \emph{scellée} au sens du protocole~: balise
git \code{deff25c}, 82 fichiers C\#, 14\,891 lignes de code, 228 tests
unitaires, trois ans de développement intermittent par le développeur et
ses prédécesseurs d’agents. Le candidat cumule les attributs d’un cas
exigeant~: temps réel audio, interactions noyau (pilote de carte virtuelle,
installateur), multi-plateforme, et un patrimoine de défauts d’installation
documentés.

L’étude poursuit trois questions de recherche. \textbf{RQ1}~: la discipline
des trois couches est-elle applicable de bout en bout à une application
réelle complète, sans exception ni accommodement~? \textbf{RQ2}~: quel est
l’effet mesurable de cette discipline sur la défense (tests, avertissements)
et la structure (taille des unités, testabilité) du code~?
\textbf{RQ3}~: que révèle la vérification de l’oracle que le
développement classique ne révèle pas~? Le design est une \emph{étude de
cas unique} au sens de Yin \cite{yin2018}, conduite et rapportée selon les
recommandations de Runeson et H{\"o}st \cite{runeson2009}~; l’unité
d’analyse est le recodage complet, de la base scellée au cristal final. Les
limites inhérentes à ce design — absence de condition témoin, validité
externe — sont traitées en section~\ref{sec:menaces} plutôt que tues.

L’instrumentation est triple. Premièrement, une intégration continue
exécutant construction et tests sur trois systèmes d’exploitation (Linux,
Windows, macOS), activée à partir de la cristallisation finale.
Deuxièmement, des empreintes cryptographiques des données embarquées
(bibliothèques de profils), reproductibles par quiconque exécute la
commande de hachage sur les deux dépôts. Troisièmement, un audit de parité
conduisant famille par famille — les quatorze familles de fonctionnalités
du fichier de lecture de la base scellée — avec, pour chacune, les modules
correspondants, la méthode de vérification et le verdict. Aucune de ces
mesures n’a été conçue après coup~: elles procèdent de la règle R5 elle-même,
appliquée à chaque vague.

\subsection{Application de la méthode}\label{sec:application}

Le recodage s’est déroulé en sept vagues, chacune close par une sauvegarde
complète (dépôt git local comme source de vérité, miroir chiffré Drive avec
bundle d’historique restaurable, miroir GitHub privé) et par la vérification
complète~: construction à 0 erreur et 0 avertissement, 100~\% des tests
verts. L’instanciation des couches est directe~: la couche $\IB$ est
constituée de 44 spécifications en texte structuré, une par module, écrites
\emph{avant} le code~; la couche $\BB$ est le code source versionné, scellé
par balises~; la couche $\CB$ est la carte du dépôt — un index à adresses
stables \code{M-<DOM>-NN}, un statut par module~; le journal $\DB$ du dépôt
compte quinze décisions, chacune avec ses alternatives rejetées. L’oracle
est la suite de tests \emph{portée puis étendue}~: la suite 2.2.2 encode le
comportement réellement validé du terrain, et chaque constante numérique de
référence (tolérances, amplitudes, fréquences) en a été extraite — la
décision D5 du dépôt pose explicitement que les tests sont le contrat
exécutable de la couche I.

Le tableau~\ref{tab:regles} (section~\ref{sec:regles}) a servi de table de
correspondance normative tout au long des vagues~: chaque règle du
protocole y trouvait son homologue concret dans le dépôt. Deux incidents
méritent d’être consignés car ils éprouvent la gouvernance plutôt que le
code. D’une part, l’environnement d’exécution a été \emph{réinitialisé} en
cours de projet — répertoires de travail et journal effacés~; la
restauration complète depuis les sauvegardes n’a coûté aucune perte de
travail, validant en conditions réelles la redondance exigée par la
gouvernance. D’autre part, une demande du développeur en fin de
démonstration — fusionner les modules et réécrire « en natif machine » —
a été refusée par décision motivée D15 après mesure~: l’analyse a établi
que la fusion détruirait les adresses stables et les spécifications par
responsabilité sans aucun gain d’exécution, le fil audio tourne déjà à
152~$\times$ le temps réel avec zéro allocation par bloc
(section~\ref{sec:resultats}). Ce refus est l’illustration la plus directe
de la fonction objectif lexicographique (définition~\ref{def:lex})~: la
demande promettait un gain d’\emph{optimisation} au prix d’une perte de
\emph{sécurité} de gouvernance — inadmissible quel que soit le gain.

\subsection{Observations qualitatives~: trois comparatifs
avant/après}\label{sec:comparatifs}

\subsubsection{Le seuil magique}\label{sec:seuil}

L’analyseur de balayage de la base scellée rejetait les captures invalides
au-dessus d’un seuil pic/médiane de~8 — un littéral sans nom, dupliqué dans
le message d’erreur, sans justification ni test. Le défaut était réel~: la
déconvolution divise par le spectre du balayage au carré, ce qui agit
comme un différenciateur hors bande~; une capture de \emph{bruit seul}
produit un faux pic à rapport mesuré d’environ 186 — largement au-dessus
de~8. Concrètement, un utilisateur ayant mal routé sa capture se voyait
attribuer une « correction » aléatoire au lieu du message d’aide. Le
protocole ne s’est pas contenté de corriger~: il a d’abord établi les faits
par diagnostic empirique (bruit blanc pur~: rapport 186~; bruit passe-bas
réaliste~: 9,2~; balayage réel même atténué~: au-delà de $10^{5}$, mesuré à
$1{,}2 \times 10^{7}$), puis consigné la décision D8 — un seuil de 1000
rejette le bruit seul avec quatre décades de marge et conserve quatre
décades de marge aux balayages faibles.

Le listing~\ref{lst:seuil} montre la conversion. Le défaut est désormais
verrouillé à quatre niveaux~: constante nommée dans le code, justification
documentée, contrat dans la spécification I du module, test exécutable qui
rejette le bruit seul.

\begin{lstlisting}[float=t, caption={Le seuil de détection, avant (haut) et
après (bas). Extraits abrégés des deux dépôts.}, label={lst:seuil}]
// AVANT — v2.2.2 : seuil littéral, sans nom, sans test :
if (peakMag < Math.Max(median * 8.0, 1e-9)) {
    throw new InvalidOperationException(
        "Balayage non détecté ... rapport ... < 8 ...");
}
// APRÈS — v3 : constante nommée, rationale, message dérivé :
/// Seuil pic/médiane. v2.2.2 : 8 laissait passer le bruit
/// seul (ratio ~186, effet différenciateur du 1/|X|^2) ;
/// balayage réel mesuré ~1,2e7 -> seuil 1000 : quatre
/// décades de marge de chaque côté.
public const double DetectionSnr = 1000.0;
if (peakMag < Math.Max(median * DetectionSnr, 1e-9)) {
    ... // message dérivé de la constante
}
\end{lstlisting}

\subsubsection{Le monolithe}\label{sec:monolithe}

Le modèle-vue de la table de mixage concentrait cinq responsabilités —
routage de la carte virtuelle, synchronisation des volumes, installation du
pilote, persistance, politique de préréglages — entrelacées sur 1\,075
lignes~; 39 propriétés observables et commandes y côtoyaient les
dictionnaires de routage et les minuteries. Aucun test ne peut s’approcher
d’une telle classe sans périphérique audio réel. La spécification I de
l’onglet mixage exigeait la décomposition \emph{avant} que le code
n’existe~; le recodage l’a réalisée en cinq fichiers à responsabilité
unique (399~+ 224~+ 209~+ 160~+ 138 lignes) et extrait la politique de
préréglages vers une classe dédiée derrière deux interfaces étroites.

Le gain n’est pas cosmétique. La décomposition a rendus visibles — et
donc corrigeables — deux défauts fonctionnels que le monolithe masquait~:
le message d’aide ASIO n’était jamais affiché sans pilote (la propriété
concernée restait nulle, aucun changement n’était notifié)~; le garde-fou
de liaison d’un préréglage à une sortie laissait passer le préréglage plat
par défaut, déclenchant une demande de nom puis un refus en cascade. Les
deux correctifs sont couverts par des tests, et la politique extraite est
vérifiée \emph{sans périphérique} — ce qui était impensable sur le
monolithe est devenu une suite de tests ordinaire. C’est la règle R3 — le
balisage — appliquée à l’architecture~: un module, une responsabilité, une
adresse stable dans la carte.

\subsubsection{Le caractère invisible}\label{sec:caractere}

Les magasins de réglages sécurisent les noms de fichiers avant d’écrire sur
disque. La base scellée déléguait cette sécurité au système hôte~: la liste
renvoyée par \code{Path.GetInvalidFileNameChars()} est celle de l’OS qui
exécute le code. Le défaut est invisible sur la machine de
développement Windows — et c’est précisément le problème~: un nom sécurisé
sous Linux, où les deux-points sont légaux, peut produire un fichier
invalide s’il est copié sous Windows. Il n’a été révélé que parce que la
vérification R5 de la démonstration exécute les tests sous Linux~:
l’oracle a fait son travail. Le correctif fige l’union des caractères
interdits sur au moins une plateforme cible, indépendamment de l’hôte~; le
contrat « un nom sécurisé est portable » est désormais écrit dans
la spécification I et dans le code. Trois leçons générales s’en
dégagent~: la vérification doit s’exécuter dans l’environnement qui révèle
le défaut~; un défaut de portabilité est un défaut de contrat~; et la
taille du correctif (six lignes remplacées par quinze) est sans rapport
avec son changement de classe — d’un comportement dépendant de l’hôte à
une garantie multi-plateformes par construction.

\subsection{Résultats quantitatifs}\label{sec:resultats}

Le tableau~\ref{tab:mesures} confronte les deux révisions sur les mesures
structurelles, défensives et de gouvernance. Le volume de code est
volontairement resté comparable — le recodage n’a pas gonflé
l’application, il a déplacé l’effort vers la spécification et la
vérification.

\begin{table}[t]
\centering\small
\caption{Confrontation des deux révisions.}
\label{tab:mesures}
\begin{tabularx}{\columnwidth}{@{}Y l l@{}}
\toprule
\textbf{Mesure} & \textbf{v2.2.2} & \textbf{v3} \\
\midrule
Fichiers C\# (\code{src/}) & 82 & 93 \\
Lignes de code (\code{src/}) & 14\,891 & 15\,438 \\
Lignes de tests & 4\,032 & 5\,963 \\
Tests unitaires & 228 & \textbf{437} \\
Plus gros fichier & 1\,075 l. & \textbf{497} l. \\
Plus gros modèle-vue & 1\,075 l. & \textbf{399} l. \\
Avertissements construction & tolérés & \textbf{0 / 0} \\
Spécifications d’intention & 0 & \textbf{44} \\
Modules cartographiés & 0 & \textbf{42} \\
Décisions journalisées & aucune & \textbf{D1–D15} \\
\bottomrule
\end{tabularx}
\end{table}

La suite de tests mérite le détail, car elle est le cœur de la
vérification. Les 228 tests de la base ont été portés intégralement —
c’était l’oracle — puis étendus~: la couche cœur passe de 150 à 220 tests
(parité des données embarquées, comportements aux limites, diagnostics
injectables), la couche audio de 78 à 142 (fournisseurs, bus principal,
anti-larsen vérifiable sous Linux, rognage exact), et une suite
applicative de 75 tests apparaît, rendue possible par la décomposition des
modèles-vues. La non-régression se démontre famille par famille~: les
quatorze familles de fonctionnalités du fichier de lecture de la base
scellée ont été auditées point par point — parité intégrale, 14/14. Les
données embarquées sont contrôlées par empreintes~: la bibliothèque de
corrections casques (321 modèles, 58 marques) et les profils de mesure
micro (36 entrées) sont identiques octet pour octet à la base scellée,
avec des sommes md5 reproductibles sur les deux dépôts. Trois écarts
volontaires subsistent — le seuil D8 (défaut connu corrigé par décision
journalisée), des correctifs de défauts de la base chacun couvert par un
test, et la décomposition structurelle sans changement de comportement
observable — et leur existence même est une propriété du protocole~:
aucun n’est silencieux.

La séquence de \emph{publication} du cristal final éprouve le mécanisme
d’échec honnête du protocole. La balise v3.0.0 a échoué en intégration
continue~: trois tests relisaient un journal vivant avec un verrouillage de
fichier incompatibles avec l’écriture sous Windows — défaut invisible sous
Linux et macOS, hérité d’une connaissance de portage perdue entre les
versions. Le correctif v3.0.1, limité aux tests, a rendu la vérification
verte sur les trois systèmes mais échoué à l’empaquetage (un outil absent
des exécutants de la CI, et une archive source trop inclusive). La version
v3.0.2 a porté l’outillage d’empaquetage sur un socle portable et ajouté
un reproducteur permanent exécuté \emph{avant} toute publication. Trois
balises, aucune déplacée~: chaque échec a été corrigé en patch à contrat
stable, conformément à la gouvernance — premier exemple complet du
mécanisme « échec de CI $\to$ patch honnête » que la règle R5 rend
obligatoire.

Enfin, la performance a été mesurée après une demande d’optimisation
« maximale » (décision D15 du dépôt)~: sur une machine à 2~vCPU partagés,
à 48~kHz stéréo par blocs de 10~ms, la chaîne d’égalisation à 30 bandes et
2 canaux traite un bloc en 66~µs — 152~$\times$ le temps réel — avec
\emph{zéro} octet alloué sur 100\,000 blocs consécutifs (1\,000~s de son)~;
le tampon circulaire producteur-consommateur coûte 0,18~µs par bloc~; la
régénération d’un balayage de 12~s prend 10,7~ms. Il n’existait aucun
problème de performance à corriger — conclusion qui fonde le refus motivé
de la réécriture « native machine » évoqué en
section~\ref{sec:application}.

\subsection{Reverse-specs de l’installateur scellé}\label{sec:reverse}

L’histoire du produit est marquée par trois défauts d’installateur
consécutifs (versions 2.2.0 à 2.2.2), dont l’analyse des journaux et des
commits révèle la classe commune~: des contrats non écrits, dupliqués en
plusieurs exemplaires. Plutôt que de recoder l’installateur (ce qui
ajouterait un chantier non validé à un chantier en attente de validation),
la démonstration a écrit ses \emph{spécifications rétro-actives}
(\emph{reverse-specs})~: huit modules M-INS rédigés en relisant ligne à
ligne le code source scellé. Le binaire reste la référence
exploitable~; le contrat, lui, devient diffable, versionné et testable. La
règle R2 — l’intention d’abord — s’applique à l’envers du temps~: le
contrat précède même la décision de recoder.

\begin{table}[t]
\centering\small
\caption{Les trois défauts historiques de l’installateur et les verrous
posés par les reverse-specs.}
\label{tab:reverse}
\begin{tabularx}{\columnwidth}{@{}Y Y@{}}
\toprule
\textbf{Défaut historique} & \textbf{Verrou posé} \\
\midrule
Structure Win32 déclarée de mémoire — 424 octets au lieu de 32 sur x64,
échec \code{0x6F8} (2 copies) & invariant de taille testé, contrat complet
SetupAPI / newdev / cfgmgr32 \\
Signature INF erronée (\code{\$Windows NT\$} au lieu de
\code{\$CHICAGO\$}), échec \code{0xE0000205} & validateur bloquant à la
construction et à l’installation \\
Analyse de la sortie pnputil sur libellés localisés français (3 copies) &
fragments indépendants des libellés, décodage CP850, repli par balayage,
fixtures partagées \\
\bottomrule
\end{tabularx}
\end{table}

Formellement, la reverse-spec est une \emph{extraction de spécification
sonore}~: étant donné un artefact scellé $\BB_{\mathrm{scellé}}$, on
construit $\IB'$ telle que le comportement observable de l’artefact
satisfait $\IB'$ ($\mathrm{Obs}(\BB_{\mathrm{scellé}}) \models \IB'$),
sans exiger la couverture exhaustive des comportements non observés. La
technique étend la couche I au patrimoine existant — situations que la
spécification du protocole n’avait pas prévues — et fige la connaissance~
: les spécifications dérivent du code relu, pas d’une mémoire de session.
Si la vague de recodage de l’installateur est un jour lancée, elle partira
de contrats déjà testés~; si elle ne l’est jamais, le comportement validé
reste spécifié pour quiconque viendra après.

% =====================================================================
%  5. DISCUSSION
% =====================================================================
\section{Discussion}\label{sec:discussion}

\subsection{Ce que l’étude établit}

Un protocole se juge à ce qu’il rend impossible. Sur les trois comparatifs
comme sur l’extension rétro-active, l’étude établit que la classe de
défauts « contrat non écrit » — valeur-seuil enfouie, structure déclarée de
mémoire, comportement dépendant de l’hôte, analyse de libellés localisés —
est systématiquement convertie en contrat nommé, documenté, testé et
journalisé. Les règles ont opéré sans exception sur un périmètre réel et
complet~: 42 modules cartographiés, 44 spécifications écrites avant le
code, dix commits tous précédés de la vérification complète, sept vagues
closes par une sauvegarde systématique. L’oracle a joué son rôle jusque
dans la révélation de défauts invisibles sur la machine de développement
(le caractère interdit dépendant de l’hôte,
section~\ref{sec:caractere}~; le verrouillage de journal révélé par la CI
Windows, section~\ref{sec:resultats}). La parité est démontrable par des
tiers — empreintes reproductibles, quatorze familles auditées, écarts
documentés — et la structure s’est améliorée sans gonflement~: à volume de
code comparable, la défense double (228~$\to$~437 tests, avertissements
tolérés~$\to$~zéro) et le plus gros fichier applicatif passe de 1\,075 à
399~lignes.

La démonstration valide aussi, indirectement, la \emph{gouvernance}. Le
journal des décisions — quinze entrées en cinq sessions, chacune avec ses
alternatives rejetées — a servi de mémoire d’arbitrage effectivement
consultée~: aucune révision silencieuse, aucun accommodement improvisé, y
compris sous la pression d’une réinitialisation d’environnement en cours
de projet. L’anti-enfermement n’est pas théorique~: l’historique, les
contrats et l’état sont récupérables par un agent neuf depuis les seuls
artefacts du dépôt. Et la séquence v3.0.0~$\to$~v3.0.2 (deux échecs de CI
honorés puis corrigés en patchs à contrat stable) montre que la
discipline tient \emph{dans l’échec}, ce qui est sa véritable épreuve.

\subsection{Implications pour le génie logiciel assisté par IA}

Pour l’AISE, l’étude suggère un déplacement du regard~: à capacité du
modèle constante, la qualité du résultat agent-first dépend du
\emph{processus} qui encadre l’agent. Les modes d’échec documentés de la
littérature — code plausible mais inexact, dérive silencieuse hors
périmètre, dépendances hallucinées \cite{hou2024} — ne sont pas que des
propriétés des modèles~; ce sont des propriétés de \emph{couples}
modèle-processus, et le second terme est actionnable dès aujourd’hui. Les
trois mécanismes que l’étude a éprouvés se transposent à tout
outillage~: (i)~la spécification écrite avant le code, qui force
l’intention à précéder l’exécution et rend chaque évolution
\emph{explicable}~; (ii)~le vérifie-puis-écrit (R5), qui transforme la
confiance dans le modèle en vérification de l’artefact — l’agent peut
halluciner, le commit reste propre~; (iii)~la cristallisation à
pré-versions explicites, qui distingue « vérifié par l’oracle » de
« validé sur le terrain » et interdit la confusion des niveaux de preuve.
La reverse-spec, enfin, adresse le problème de l’\emph{embarquement} de
patrimoine~: face à un legs sans documentation, l’agent ne repart pas de
zéro — il extrait d’abord le contrat du code relu, puis le fait valider.

\subsection{Coûts et frictions observés}

L’honnêteté impose de chiffrer la contrepartie. La discipline a un coût
d’écriture~: 44 spécifications et un journal de quinze décisions
représentent un investissement documentaire significatif, supérieur à ce
qu’un développement humain classique consent — l’étude de cas l’a payé
parce qu’il s’agissait précisément de valider la discipline. Elle a un
coût de tempo~: aucune écriture n’est permanente sans vérification
complète, ce qui rend les itérations plus lentes qu’en édition libre — le
prix payé pour que la régression soit impossible \emph{par construction}
plutôt que rare en pratique. Elle a enfin un coût d’infrastructure~:
intégration continue multi-OS, triple sauvegarde, reproducteur
d’empaquetage. L’argument économique du protocole est que ces coûts sont
des \emph{investissements anti-dette}~: dans l’étude de cas, le diagnostic
complet du seuil D8 — mesure, décision, spécification, tests — a coûté une
fraction de ce qu’aurait coûté le seul dépannage du défaut en production~;
et le défaut existait depuis trois ans parce que rien ne l’avait rendu
visible. Ces observations restent qualitatives~; aucune mesure
comparative du coût total n’a été tentée (section~\ref{sec:menaces}).

\subsection{Enseignements pour le protocole lui-même}

L’étude renvoie au protocole trois enseignements concrets, que la
prochaine révision pourra instruire. Premièrement, la notion de
\emph{pré-version de cristallisation} (\code{-dev.N})~: la démonstration a
distingué « vérifié par l’oracle » de « validé sur le terrain », nuance
que la spécification ne prévoyait pas. Deuxièmement, la \emph{reverse-spec}
elle-même, qui étend la couche I aux binaires scellés hérités. Troisièmement,
la valeur du \emph{partage de fixtures} entre consommateurs d’un même
contrat externe — les spécifications de l’installateur et l’application
partagent contractuellement leurs jeux de test d’analyse de sortie, ce qui
élimine par construction la duplication qui laisse vivre les défauts en
trois exemplaires. Chacun de ces enseignements est entré au journal des
décisions du dépôt~; aucun n’a été appliqué en silence — la gouvernance
s’applique à elle-même.

% =====================================================================
%  6. MENACES SUR LA VALIDITÉ
% =====================================================================
\section{Menaces sur la validité}\label{sec:menaces}

Nous suivons la classification de Wohlin et al. \cite{wohlin2012},
complétée par les recommandations de reporting des études de cas de
Runeson et H{\"o}st \cite{runeson2009}. Aucune de ces menaces n’est
résolue~; toutes sont déclarées.

\subsection{Validité de construction}\label{sec:construction}

Les métriques mobilisées — nombre de modules, de tests, d’avertissements,
taille du plus gros fichier — sont des \emph{indicateurs} de structure et
de défense, pas des mesures directes de qualité. Le compte de tests ne
mesure ni la couverture ni la force des tests~: \emph{aucune mesure de
couverture de code n’a été collectée} durant l’étude, et l’appréciation
qualitative des extensions (« parité des données embarquées »,
« comportements aux limites ») est celle de l’agent auteur. Le compte de
modules dépend de la discipline de cartographie elle-même — l’instrument
et l’objet mesuré sont partiellement confondus. Enfin, la « parité
fonctionnelle 14/14 » est une audit de familles de fonctionnalités guidé
par le fichier de lecture de la base, pas une équivalence formelle~; elle
ne couvre ni les comportements non documentés, ni les caractéristiques non
fonctionnelles au-delà des tests.

\subsection{Validité interne}\label{sec:interne}

Il n’existe \emph{pas de condition témoin}~: aucun recodage parallèle de la
même base sans la discipline TRIADE n’a été conduit, et l’attribution des
améliorations au protocole plutôt qu’à l’attention renouvelée portée au
code est un jugement, non une mesure. La structure de l’étude est
auto-référentielle~: l’agent qui a écrit les spécifications a aussi écrit
le code, les tests, l’audit et le présent article — un biais
d’auto-confirmation que seule une évaluation indépendante pourrait
lever. Les trois comparatifs avant/après ont été \emph{sélectionnés}
comme les plus illustratifs~; les vagues ont aussi produit des
conversions moins spectaculaires, non rapportées. Les mesures de
performance ont été prises sur une infrastructure partagée à 2~vCPU, dans
des conditions favorables au meilleur cas (meilleur de cinq essais).

\subsection{Validité externe}\label{sec:externe}

C’est la menace majeure~: l’étude est \emph{monocasale} (N~$=$~1). Un
projet, un domaine (audio temps réel), une pile technologique
(C\#/.NET~8/Avalonia), un agent (un unique modèle, dans un unique
environnement), un développeur, et — pour le protocole lui-même — ses
propres auteurs, puisque la démonstration a servi à valider la méthode
qui l’encadrait. Aucun claim de généralisation n’est avancé~: les
mécanismes rapportés sont \emph{plausibles} dans d’autres contextes, pas
\emph{établis}. La réplication sur d’autres projets, piles et agents,
conduite par des tiers, est la condition de toute montée en généralité
\cite{runeson2009}. De même, l’applicabilité à des équipes humaines
nombreuses — dont les incitations, contrairement à celles d’un agent,
comprennent la vitesse de livraison individuelle — est une question
ouverte que le protocole n’a pas encore affrontée.

\subsection{Validité de conclusion}\label{sec:conclusion-val}

Les affirmations de performance reposent sur des micro-benchmarks en
meilleur cas sur infrastructure partagée~; elles établissent l’\emph{absence
d’un problème de performance à corriger}, pas des performances
représentatives de déploiements réels. Le « zéro régression » est borné
par l’étendue de la suite de tests et de l’audit de familles~; il ne
préjuge pas des comportements non couverts — l’épisode du verrouillage de
journal sous Windows (invisible sous Linux/macOS) rappelle qu’une suite
verte est une \emph{preuve conditionnelle}, conditionnelle à
l’environnement d’exécution. L’épisode v3.0.0~$\to$~v3.0.2, présenté comme
validation du mécanisme d’échec honnête, est un événement unique~; sa
généralité n’est pas démontrée.

\subsection{L’écart conceptuel~: refactorisation source ou génération
binaire native~?}\label{sec:ecart}

Une menace spécifique mérite un traitement distinct, car elle touche au
rapport entre la démonstration et l’objet annoncé du protocole. La
spécification TRIADE v0.3 promet, à terme, un développement \emph{direct
en code machine natif scellé} — la couche B comme suite d’octets écrits
via une notation assemblée et vérifiée par round-trip, le compilateur
disparaissant du pipeline. L’étude de cas n’a rien démontré de tel~: la
couche B y est instanciée comme \emph{code source managé}, compilé par une
chaîne standard (compilation .NET, exécution avec compilation à la volée),
et les mécanismes spécifiques au natif — notation, round-trip, verrou
croisé cryptographique par empreintes d’octets, patch à chaud par
trampoline, adresses garanties par remplissage — n’ont pas été exercés.

L’étude valide donc la \emph{couche de gouvernance} du protocole — les
trois artefacts, les dix règles, la boucle, la vérification-puis-écriture,
la cristallisation, l’arbitrage lexicographique — appliquée à une
instanciation source de la couche B. Elle ne valide ni la forme de
l’artefact, ni les mécanismes natifs, ni l’argument de performance qui
les motive~; ces points restent la charge des jalons M0 à M7, non
encore franchis. Symétriquement, la motivation \emph{agent-first} du
protocole — l’économie de contexte par édition ciblée via la carte —
repose sur des budgets cibles (40 / 3\,000 / 6\,000 tokens) qui sont des
\emph{objectifs de conception}, pas des mesures~: aucune instrumentation
du coût réel en tokens n’a été conduite durant l’étude. Nous posons ces
limites d’autant plus clairement que la tentation symétrique existe~:
conclure de l’étude à la validité du programme natif serait aussi
infondé que conclure de son cadre à l’inutilité de la gouvernance. Ce
que l’étude établit tient en une phrase~: \emph{la discipline de
gouvernance est applicable de bout en bout, à coût documenté, à une
application réelle complète, et les défauts qu’elle a convertis en
contrats sont ceux-là mêmes qui survivaient depuis trois ans à la
maintenance classique}.

% =====================================================================
%  7. CONCLUSION
% =====================================================================
\section{Conclusion et perspectives}\label{sec:conclusion}

Nous avons présenté TRIADE, un protocole de gouvernance du développement
logiciel pensé pour des exécutants agentiques~: trois artefacts versionnés
(intention, binaire, carte), dix règles à sanction mécanique, un verrou
croisé cryptographique, une fonction objectif lexicographique, et une
boucle de mise à jour dont le commit est gardé par la vérification.
L’article en a donné une présentation formalisée et l’a confrontée à sa
première validation à l’échelle~: le recodage complet de l’application
RoomEq — 44 spécifications avant le code, 42 modules cartographiés, 437
tests contre 228, zéro avertissement, parité fonctionnelle vérifiée,
séquence de publication incluant deux échecs d’intégration continue honorés
— avec une technique au passage~: les spécifications rétro-actives de
l’installateur scellé. Les menaces sur la validité ont été traitées comme
une contribution à part entière~: cas unique sans témoin,
auto-vérification par l’agent auteur, et écart assumé entre la
refactorisation source démontrée et la génération binaire native promise.

Trois perspectives découlent directement de ces limites.
Premièrement, la \emph{réplication}~: appliquer le protocole à d’autres
projets, piles et agents, idéalement par des équipes indépendantes, pour
convertir la plausibilité en généralité. Deuxièmement, la
\emph{réalisation}~: franchir les jalons M0–M7 de la feuille de route —
chaîne notation-octets-exécutable sans compilateur, patch à chaud sur code
en exécution, portabilité inter-OS puis inter-ISA — qui seuls valideront
le cœur natif du protocole. Troisièmement, l’\emph{instrumentation}~:
mesurer ce que l’étude a seulement ciblé — coût réel en tokens par
édition, couverture de code, coût total de la discipline comparé à un
développement non gouverné. Le fil conducteur reste celui de
l’introduction~: développer avec un agent exige de l’agent qu’il prouve
ce qu’il affirme, et du projet qu’il survive à l’agent. La démonstration
RoomEq établit que la première moitié de cette exigence est tenable à
l’échelle d’une application réelle~; la seconde est, par construction, en
travail permanent.


\bibliography{refs}

\end{document}
